\section{The figure of merit: the angle of each track}\label{sec:requirement}
\sectionstate{partial}{new section (Dylan 2026-11-10: ``a diagram and section for arguing why angular resolution is most
important \ldots{} a quick toy simulation with visualization''). Toy and figure: \texttt{papers/paper1\_v2/figs/make\_pointing.py},
numbers in \texttt{figs/out/pointing.json} (keys \texttt{air\_only}, \texttt{air\_only.pair\_symmetric},
\texttt{slope\_sensitivity}). 2026-11-11: made general (Dylan: no accounting of the n\_TOF capsule or of this chamber's
window here, the argument must carry to LNL and to a thinner window): air only, generic source size. 2026-11-12
(Dylan): the source is taken as extended; the compact case kept to one paragraph. 2026-11-19 (Dylan): rebuilt from the
ground up. The section used to open on the pointing to the target; it now opens on the opening angle, shows that
position and time enter the measurement through the slope of the segment, and keeps the pointing as one paragraph.}
\intent{Before the detector is described: the search measures an opening angle, the angle between two short segments
each measured inside one chamber. Neither where a segment is nor when the event happened enters it. The space and
time measurements of the chamber matter because the slope is made of them (differences across the gap), so the
angular resolution is the one number that reaches the physics. The pointing to the target, the one use of the
position, is limited by the angle too. About 1\textdegree{} is enough because the air alone scatters more. General:
nothing here depends on the target or on the foils of this chamber.}

An X17 search comes down to the a measurement of angles. The two leptons share the energy $E$ of the transition, and with
$y=(E_+-E_-)/E$ and, neglecting the electron mass for simplicity, the invariant mass of the pair is
\begin{equation}
  m^2=E^2\,(1-y^2)\,\sin^2(\theta/2)\,,\label{eq:mass}
\end{equation}
with $\theta$ the opening angle (inset of \crefp{fig:pointing}{a}). The transition under study constrains $E$ and the calorimetry behind the trackers measures $y$. The
tracker, then, contributes $\theta$. An X17 of 16.8\,MeV appears as an edge in $\theta$, at 109\textdegree{} for the
20.6\,MeV transition of $^4$He and at 136\textdegree{} for the 18.15\,MeV transition of $^8$Be, on an IPC background
that falls smoothly with $\theta$. There, one degree of opening angle is 0.10 and 0.06\,MeV of mass. The angular
resolution therefore sets how sharp the edge stays and how much of the background lies under it.

\begin{figure}[tbp]
  \centering
  \figorbox{pointing}{to-scale side view with the pair kinematics inset, and direction error against lepton energy}
  \caption{What the angle is worth. (a) One chamber and an extended target to scale, with one lepton (black; its
  deflection in the air is drawn about five times too large for an 8\,MeV electron). The chamber measures a segment
  about 30\,mm long, about 230\,mm from the target (an idealized geometry). One degree on its direction is 0.5\,mm
  across the drift gap and 4\,mm at the target: the measured angle, extrapolated back (blue, dashed), misses the
  origin by $b$; the dotted line continues the initial direction. Inset, not to scale: the pair where it is made, with the opening
  angle $\theta$ and the energies $E_+$ and $E_-$ of the two leptons, drawn at the 109\textdegree{} of the X17 edge in
  $^4$He. (b) Error of the measured angle on the initial
  direction of one lepton, in one projection, against its kinetic energy, with 20\,cm of air as the only material.
  Gray: the scattering in the air, which is what a perfect chamber would give. Blue: a chamber of 1\textdegree{}
  (dashed line) added to it; the line is analytic and the points a toy simulation that scatters the leptons through
  the air.}
  \label{fig:pointing}
\end{figure}

Each direction is measured inside one chamber. With an extended source, such as a gas target in a neutron beam, the
point where a pair was made is not known beforehand, so the direction of a lepton cannot be taken from a line drawn
between the target and the chamber. It has to come from what the chamber sees: a straight segment about 30\,mm long,
about 20\,cm from the target (\crefp{fig:pointing}{a}). The opening angle is the angle between two such segments.
Neither the position of a segment nor the time of the event enters it: a segment moved sideways, or an event moved in
time, gives the same $\theta$. Its error is that of the two slopes, $\sqrt2\,\sigma_\theta$ for the component in
the plane of the pair, with $\sigma_\theta$ the angular resolution of a chamber.

Space and time do matter, but inside the segment, where they make the slope. The ionization of an inclined track
reaches the readout displaced across the strips by $\Delta x$ between the two ends of the drift gap and spread in
arrival time by $\Delta t$ (\crefp{fig:layout}{b}), and with the drift velocity $\vd$
\begin{equation}
  \tan\theta_\mathrm{track}=\frac{\Delta x}{\vd\,\Delta t}\,.\label{eq:slope}
\end{equation}
For a track near the normal, one degree is 0.5\,mm of $\Delta x$ across a 30\,mm gap. For a track inclined by
30\textdegree{} one degree is also 1.2\,mm of depth, which is 34\,ns of drift at 36\,\textmu m/ns, and an error of
2\,\% on $\vd$ is 0.5\textdegree. These are differences along one track in one chamber: an offset common to the segment,
the position of the chamber or the time of the trigger, cancels in the slope. The strip pitch, the diffusion, the
time sampling and the calibration of the drift velocity are the ingredients of the angle, and $\sigma_\theta$ is the
one number through which they reach the mass.

The position of the segment enters a second, lesser quantity: the extrapolation of the track back to the target,
which tells whether two tracks come from one point and whether that point lies in the target. This rejects
background; it does not measure the mass. With the position known to $\sigma_x$ at a distance $d$ from the target,
the track misses its origin by
\begin{equation}
  \sigma_b=\sqrt{\sigma_x^2+(d\,\sigma_\theta)^2}\,.\label{eq:pointing}
\end{equation}
With $d=233$\,mm from the target to the readout plane, an angle error of 1\textdegree{} is 4.1\,mm at the target.
A position error of 0.5\,mm is 1.5\,\% of the variance beside it, and the position would have to be as poor as 4\,mm
to matter as much as one degree: the lever arm is eight times the length of the segment. The event time counts for
even less. It places the segment along the drift, where 40\,ns is 1.5\,mm, and otherwise serves to associate the two
leptons with each other and with the scintillators. A toy simulation that throws tracks from the target over the
acceptance of a chamber, smears the measured position and slope and extrapolates them back reproduces
\cref{eq:pointing}, about 10\,\% above it because inclined tracks have a longer lever arm. So the one place where
the position is used is limited by the angle as well.

What sets the useful scale for $\sigma_\theta$ is scattering before the measurement, and the air between the target
and the chamber is enough to set it. Over 20\,cm of air an electron is deflected by 3.2\textdegree{} rms in each
projection at 4\,MeV, by 1.7\textdegree{} at 8\,MeV and by 0.9\textdegree{} at 16\,MeV
(\crefp{fig:pointing}{b}). That is more than 1\textdegree{} below 14\,MeV, before any target wall, window or foil is
counted (those of the present chambers are in \cref{sec:material}). A chamber much better than 1\textdegree{} would
therefore gain little, and a worse one would cost: for a symmetric pair at the $^4$He edge, two leptons of 9.8\,MeV,
the opening angle is known to 2.0\textdegree{} with perfect chambers, to 2.4\textdegree{} with chambers of
1\textdegree{} and to 3.5\textdegree{} with chambers of 2\textdegree, which is 0.21, 0.25 and 0.36\,MeV of mass.

If the source is compact instead, a beam spot of a few millimeters on a thin target, the line from the known origin
to the measured position is the better estimate of the direction (for a source below about 7\,mm rms at 8\,MeV),
and the measured angle is left to confirm that the track comes from the source. Even then the position is not
demanding: 1\,mm at this distance is 0.25\textdegree.

A tracker for these leptons therefore needs an angular resolution of about a degree, which means following the
ionization of a track to a fraction of a millimeter across its drift gap, in the strip coordinate and in depth. The
position of the track as a whole is needed only to a few millimeters and the time of the event only for the
coincidence, and there should be as little material as possible in front of the drift volume. \Cref{sec:tracking}
gives what the chambers deliver against this.
